% Chapter VI component. Standard LaTeX + amsmath, amssymb only.
% Self-contained finite-dimensional proofs; no external theorem package required.
\section{Restoration of a local state and retention of a physical history}
\label{sec:restoration-history}

The source interface is
\begin{equation}
 W8+X1\longrightarrow W9\longrightarrow W8'.
 \label{eq:write-interface}
\end{equation}
Here $W8$ names the resting, non-contact local substrate state; $X1$
names the supplied interaction; and $W9$ names the temporary write state.
The prime marks return after interaction, not a new substrate atom.
The source's intended distinction is that a local state can return while
an event remains recorded elsewhere. This section gives that distinction
an explicit quantum-information comparator. It does not identify the
symbols $W8$, $X1$, and $W9$ with Hilbert-space dimensions, derive their
exact SAM/SLC grammar, or prove a substrate microdynamics.

Chapter II already distinguishes an address from the history of its use.
Returning to a retained address does not by itself say whether a record
survives. Nor does a finite alphabet supply unlimited storage: a physical
history requires distinguishable states of some specified recipient.
The same discipline applies here. We must name the restored subsystem,
the other degrees of freedom, and the operations by which a record can
be read.

\subsection{What a unitary reset can and cannot do}
\label{subsec:reset-theorem}

Let $L$ be a local record cell and let $E$ include every other degree of
freedom participating in a closed finite-dimensional evolution. Suppose
$\{|m_i\rangle_L\}_{i=1}^n$ are orthonormal pointer states and $E$ starts
in a fixed pure state $|e_0\rangle$. An exact reset of $L$ has the form
\begin{equation}
 U\bigl(|m_i\rangle_L|e_0\rangle_E\bigr)
   =|0\rangle_L|e_i\rangle_E.
 \label{eq:unitary-reset}
\end{equation}
\noindent\textbf{Reset--recipient proposition.}
If $U$ is unitary, then $\langle e_i|e_j\rangle=\delta_{ij}$.
Consequently $E$ cannot also return to one fixed state when $n>1$.
The recipient needs at least $n$ mutually orthogonal states.

\noindent\emph{Proof.}
Preservation of inner products gives
\begin{equation}
 \delta_{ij}
 =\langle m_i,e_0|m_j,e_0\rangle
 =\langle 0,e_i|0,e_j\rangle
 =\langle e_i|e_j\rangle.
\end{equation}
An exact pure local output factors from the rest, so this argument also
covers a purported reset whose output was initially written as a general
joint pure state. A mixed initial environment may be purified by adding
its reference to $E$; the conclusion then concerns that enlarged
recipient. It does not guarantee a record readable from an arbitrarily
chosen proper part of $E$. $\square$

There is an important scope condition. If a separate system $S$ already
contains a correlated copy, the input is not the fixed-environment input
of Eq.~\eqref{eq:unitary-reset}. In the more general pure-state mapping
\begin{equation}
 |s_i\rangle|m_i\rangle|e_0\rangle
 \longmapsto |s_i'\rangle|0\rangle|e_i\rangle,
\end{equation}
unitarity requires, for $i\ne j$,
\begin{equation}
 \langle s_i'|s_j'\rangle\langle e_i|e_j\rangle=0.
 \label{eq:reset-factor}
\end{equation}
If $S$ returns to one common pure state, the recipient states must be
orthogonal. If $S$ instead retains orthogonal alternatives, it can itself
carry the distinction; a new record in $E$ is not compulsory. For example,
reversing a controlled write resets the cell while retaining the source
bit. The assertion that ``every reset must write a new history register''
would therefore be false. The invariant is distinguishability of the
complete closed state, not the permanence of a selected local mark.
For mixed or entangled subsystems, surviving information may be in joint
correlations rather than readable in any one marginal.

\subsection{Three registers: a write, a restoration, and a retained marker}
\label{subsec:three-register-write}

Use three qubits: route $S$, local cell $L$, and history recipient $H$.
The history register is initially blank. Define
\begin{equation}
\begin{gathered}
C|r,\ell,h\rangle=|r,\ell\mathbin{\oplus}r,h\rangle,
 \\
 T=\operatorname{SWAP}_{LH},
\end{gathered}
\end{equation}
where $r,\ell,h\in\{0,1\}$ and $\oplus$ is addition modulo two.
Both maps permute the full eight-state computational basis, so both
are unitary without a restriction to the intended input states.
Their composition is
\begin{equation}
\begin{gathered}
U=TC,\\
 U|r,\ell,h\rangle
     =|r,h,\ell\mathbin{\oplus}r\rangle.
\end{gathered}
\label{eq:three-register-unitary}
\end{equation}
For a normalized coherent route input,
\begin{align}
 |\Psi_{\rm in}\rangle
  &=(\alpha|0\rangle+\beta|1\rangle)_S|0\rangle_L|0\rangle_H,
  \nonumber\\
 |\Psi_{\rm write}\rangle
  &=\alpha|000\rangle+\beta|110\rangle,
  \label{eq:write-stage}\\
 |\Psi_{\rm restored}\rangle
  &=\alpha|000\rangle+\beta|101\rangle
  \nonumber\\
  &=|0\rangle_L\bigl(\alpha|00\rangle_{SH}
                        +\beta|11\rangle_{SH}\bigr).
 \label{eq:restored-stage}
\end{align}
The local cell is exactly restored. Conditioned on the route alternative,
the history states are $|0\rangle_H$ and $|1\rangle_H$, which can be
perfectly discriminated. For a coherent input, however, the final state
is an entangled superposition, not a derivation of one selected outcome.
The comparator exhibits a readable pointer correlation and local
restoration. It does not derive collapse or a unique realized history.

The analogy to Eq.~\eqref{eq:write-interface} is deliberately limited:
resting cell, interaction-enabled write, restored cell with a retained
recipient correlation. The interaction $X1$ is not treated as an extra
qubit or an arithmetical increment. The temporary $W9$ designation may
refer to the interaction configuration; it need not mean that the cell
alone has changed in every branch. In particular, the $r=0$ branch leaves
its pointer value at zero even during the controlled write.

The inverse exists on every input: $U^{-1}=C T$. With joint control of
$S,L,H$, it returns Eq.~\eqref{eq:restored-stage} to
$|\Psi_{\rm in}\rangle$. The phase information was never deleted from
the complete state. By contrast, tracing out the record gives
\begin{equation}
 \rho_S=|\alpha|^2|0\rangle\langle0|
          +|\beta|^2|1\rangle\langle1|
 \label{eq:restored-route-marginal}
\end{equation}
both after the write and after the restoration. Thus the swap does not
restore route-only interference. Global coherence is preserved throughout;
its observability depends on access to the correlated degrees of freedom.
Joint reversal, or suitable measurements of the record with conditioned
route statistics, can reveal that coherence. An operation on the
inaccessible record alone cannot change the unconditioned route marginal.

For example, take $\alpha=1/\sqrt2$ and $\beta=e^{\mathrm{i}\phi}/\sqrt2$.
The joint observable $X_S X_H$ has expectation $\cos\phi$ after restoration,
whereas every route-only expectation is independent of $\phi$.
This is the precise distinction between preservation in the closed system
and accessibility in a restricted domain.

This device stores one binary marker. It does not store an arbitrary
sequence of writes. If $H$ is already occupied, the swap returns that
previous content to $L$, rather than resetting $L$ to zero. Repeated
restoration needs additional blank capacity, a deliberate reversible
compression when possible, or an explicitly modeled disposal channel.
Perfectly distinguishing $N$ possible complete histories requires $N$
orthogonal global record states and hence a record-space dimension at
least $N$. This is a storage statement, not a heat formula.

\subsection{Pointer copying is not universal quantum cloning}

For a declared orthogonal pointer basis, a controlled operation can copy
its label:
\begin{equation}
 |r\rangle|0\rangle\longmapsto |r\rangle|r\rangle.
\end{equation}
Applied to a superposition, it creates the correlated state in
Eq.~\eqref{eq:write-stage}; it does not produce two copies of
$\alpha|0\rangle+\beta|1\rangle$. Indeed, let the overlap of two normalized inputs be $z$ and allow the
cloner to finish in input-dependent auxiliary states $|a_\psi\rangle$ and
$|a_\varphi\rangle$. Inner-product preservation would require
$z=z^2\langle a_\psi|a_\varphi\rangle$. For $0<|z|<1$, this implies
$1\le |z|$, a contradiction. Universal cloning is therefore impossible
even when the auxiliary state is allowed to change.

Transfer is different. A swap implements
\begin{equation}
 |\psi\rangle_L|0\rangle_H\longmapsto
 |0\rangle_L|\psi\rangle_H
\end{equation}
for every $|\psi\rangle$, preserving its entanglement with any reference.
The input remains present once, at its new location. Resetting $L$
while leaving both the recipient and all other degrees of freedom fixed
would instead delete an unknown input, contrary to inner-product
preservation. These elementary statements suffice for this interface;
no universal copying or unexplained destruction of quantum information is being assumed. Related no-deleting and no-hiding results have their own hypotheses; no claim that information must reside only in correlations is used here.\cite{PatiBraunstein2000,BraunsteinPati2007}

\subsection{A physical record needs an accessible decoder}
\label{subsec:accessible-decoder}

The decoder criterion of Eq.~\eqref{eq:operational-record} applies to the nominated recipient, not merely the restored cell. For equal priors, the three-register example has trace distance one between its conditional history states and zero between its conditional local-cell states. The history permits perfect route decoding; the restored cell permits only chance performance. The Helstrom optimum is attainable only if its measurement belongs to the allowed set.

Durability adds a time and noise specification. A useful requirement is
that a stated decoder recover the label with error at most $\epsilon$
after a declared storage interval and channel. Orthogonality at the moment
of writing does not by itself establish durability. For example, encoding
a classical bit in an odd number $n$ of identical pointer cells and
reading by majority vote gives
\begin{equation}
\begin{gathered}
P_{\rm err}\le
 \exp\{-2n(\tfrac12-p)^2\},\\
 p<\tfrac12,
\end{gathered}
\end{equation}
if, conditional on the stored bit, the cell errors are independent and
each has probability at most $p$. For error indicators $E_j$ and $t=1/2-p>0$, independence and the
bounded-variable exponential-moment inequality give
\begin{equation}
 \Pr\!\left\{\sum_j(E_j-\mathbb E E_j)\ge nt\right\}
 \le \inf_{\lambda>0}e^{-\lambda nt+n\lambda^2/8}
 =e^{-2nt^2}.
\end{equation}
Majority failure is contained in the displayed event. Correlated failures need not satisfy it: one
common flip of all cells leaves a majority decoder with the same error
probability as a single cell. This example is an optional classical
storage mechanism, not a derived SAM durability law.

Finally, a history description is not an additional quantum observable.
If two proposed histories yield the same accessible density operator,
every accessible POVM gives the same outcome probabilities. Calling one
state $W8$ and the other $W8'$ cannot change this. If a history can be
read, its evidence must reside in distinguishable accessible states or
correlations, with a specified decoding operation. This is the physical
content of the prime: local return is compatible with a retained record,
but the notation alone neither guarantees that record nor makes it
accessible.

No thermodynamic cost is inferred from these unitary maps. A Landauer
bound belongs to a separately declared cyclic erasure protocol with a
thermal reservoir, temperature, initial logical ensemble, and work/heat
accounting. Neither the number of local addresses nor the logarithm of a
record-space dimension is, by itself, dissipated heat. The present result
is narrower: a local restoration can transfer a physical distinction;
it cannot establish the disappearance of that distinction from a closed
unitary account.

